\section{Background and Motivation}
\label{sec:background}

Source-level JML specifications are well established: the language~\cite{jml,leavens2006design} and its checker~\cite{openjml} have a thirty-year history, and the inference engine of Article~1 produces them at scale. The gap that motivates this article is downstream of inference: once a specification is known for a method, how does it travel to the consumer that needs it?

Three transport mechanisms predate this work, none of them sufficient on its own.

\paragraph{JML stub files.}
Files of the form \texttt{<class>.jml} or \texttt{<class>.refines-jml} carry signatures and clauses without bodies. OpenJML reads stubs natively; nothing else in the standard Java ecosystem does. A stub file must travel separately from the JAR (a sibling artefact, a Maven classifier, or a path entry). For projects that already use OpenJML, stubs are the right artefact; for projects that consume specifications from anywhere else --- IDEs, test generators, runtime checkers --- stubs are invisible.

\paragraph{Annotations.}
The Java annotation mechanism (JLS~\S9.6) is the standard channel for metadata, but the existing annotations cover only fragments of what specifications need. JSR-305 captured nullability (\texttt{@Nullable}, \texttt{@NotNull}) and was withdrawn before final standardisation; the surviving informal usage covers only nullability. The Checker Framework's annotations are type-system-shaped and do not generalise to \texttt{\textbackslash sum}, \texttt{\textbackslash forall}, or relational postconditions. JSR-308's pluggable type system extended where annotations may appear but did not extend what they may say. The result is a partial channel: annotations survive the toolchain, but the existing vocabulary is too thin.

\paragraph{Class-file attributes.}
JVMS~\S4.7.1 permits vendor-defined class-file attributes. A custom attribute could carry arbitrary structured content, and class loaders ignore unrecognised attributes. The cost is invisibility: standard reflection does not surface custom attributes, IDEs do not display them, ProGuard/R8 strip them by default, and javap requires \texttt{-private -v} to even mention them. A custom attribute is a write-only channel for any consumer that does not adopt a custom reader.

\paragraph{REST API descriptions.}
At service boundaries, OpenAPI~3.x and gRPC's protobuf carry a complementary class of metadata. OpenAPI permits arbitrary \texttt{x-*} extension properties; gRPC has no comparable extension point but supports custom Protobuf annotations. Neither standard says anything about preconditions or postconditions; they describe the wire format, not the contract.

\paragraph{The carrier this article proposes.}
A repeated runtime-retained annotation type (\texttt{@JmlSpec}, with one annotation per JML clause) plus an OpenAPI extension family (\texttt{x-jml-*}) covers the gap. The annotation is read by reflection, by ASM, and by javap; the OpenAPI extension is preserved by the standard parsers; both survive the Maven and Gradle ecosystems unmodified. Section~\ref{sec:format} specifies the format precisely; Section~\ref{sec:impl} describes the embedder/extractor implementation; Section~\ref{sec:empirical} reports the empirical evaluation.
